\chapter{Grundlagen}
\label{cha:Grundlagen}


\textit{Hier beschreibt jeder die theoretischen Konzepte, die für seine Aufgaben notwendig sind.}

% --- ANGUS ---
\section{Mechanische Konstruktionsprinzipien (Angus)}
\subsection{Stabilität bei mobilen Robotern}
\subsubsection{Einfluss des Schwerpunkts und der Stützradposition}
\subsubsection{Mechanik von Wechselakkusystemen}

\subsection{Kinematik von Service-Roboter-Armen}
\subsubsection{Freiheitsgrade (DOF) und Gelenktypen}
\subsubsection{Materialkunde für Gehäuse und Kopfmodule}

% --- SIMON ---

\subsubsection{Ressourcenmanagement auf Einplatinencomputern}

\subsection{Mensch-Maschine-Interaktion (HMI)}
\subsubsection{Technologien der Sprachsynthese und -erkennung}
\subsubsection{Netzwerkprotokolle zur Fernwartung (VNC)}

% --- HAI NAM ---
\section{Elektromechanische Integration und Design (Hai Nam)}
a
\subsection{Aktorik und Greifsysteme}
a
\subsubsection{Funktionsweise und Ansteuerung von Servomotoren}
a
\subsubsection{Mechanik von Greifklauen}
a
\subsection{Elektronik-Design und EMV}
a
\subsubsection{Kabelmanagement und Signalintegrität}
a
\subsubsection{Corporate Design Prinzipien im Roboterbau}
a
\section{KI TTS STT}
\subsection{Fourier-Transformation} 
\subsection{TTS} %text to speech
\subsection{STT}


%was ist das , funktionsprinzip,traingingsmethoden, fähigkeiten von kleinen LLMs 


